\chapter{Static Range Queries}\label{ch:static}

A \defn{range query} asks for an aggregate (a sum, a minimum, a gcd, the number of distinct values) over a contiguous part $[l, r)$
of an array. When the array never changes, preprocessing buys very fast queries, and when all queries are known in advance, they can
be reordered to share work. This chapter covers the static and offline toolbox: sparse tables, their disjoint variant for any
associative operation, Cartesian trees, square-root decomposition, and Mo's algorithm with its variants for updates and for trees.
Chapter~\ref{ch:segtree} adds updates to the picture.

\begin{chapterintro}{\gold\ \plat\ \ioi}{Chapter~\ref{ch:prefix} (prefix sums, monotonic stacks) and Chapter~\ref{ch:trees}
(Euler tours, LCA).}
  \item answer idempotent range queries in $\Oh{1}$ with a sparse table, and any associative one with a disjoint sparse table;
  \item build a Cartesian tree in $\Oh{n}$, and use it to turn range minima into LCAs and to count subarrays by their minimum;
  \item split an array into $\sqrt n$ blocks when nothing else fits, and rebuild periodically;
  \item reorder offline queries with Mo's algorithm, including updates and paths in trees.
\end{chapterintro}

\begin{table}[htbp]
\centering\small
\begin{tabularx}{\linewidth}{@{}>{\raggedright\arraybackslash}p{3.6cm}>{\raggedright\arraybackslash}p{2.6cm}>{\raggedright\arraybackslash}p{2.9cm}X@{}}
\toprule
structure & preprocessing & query & works for \\
\midrule
prefix sums & $\Oh{n}$ & $\Oh{1}$ & invertible operations: $+$, $\oplus$, products with inverses \\
sparse table & $\Oh{n \log n}$ & $\Oh{1}$ & idempotent operations: min, max, gcd, AND, OR \\
disjoint sparse table & $\Oh{n \log n}$ & $\Oh{1}$ & any associative operation \\
Cartesian tree & $\Oh{n}$ & via LCA & minimum and its position, ``where is $a_i$ the minimum'' \\
square-root decomposition & $\Oh{n}$ & $\Oh{\sqrt n}$ & almost anything, also with updates \\
Mo's algorithm & offline & $\Oh{\sqrt n}$ amortised & statistics maintained by adding and removing one element \\
\bottomrule
\end{tabularx}
\caption{The static range-query toolbox.}
\label{tab:static}
\end{table}

% ======================================================================================================================
\section[Sparse table]{Sparse table \hfill \gold}\label{sec:sparse}

Prefix sums (Section~\ref{sec:prefix}) answer range sums because subtraction undoes addition. A minimum cannot be undone, but it has
another property: it is \emph{idempotent}, $\min(x, x) = x$, so counting an element twice does no harm. A \defn{sparse table} stores
the minimum of every block whose length is a power of two,
\[ \mathit{st}[k][i] = \min(a_i, \dots, a_{i + 2^k - 1}) = \min\bigl(\mathit{st}[k-1][i],\ \mathit{st}[k-1][i + 2^{k-1}]\bigr), \]
$\Oh{n \log n}$ values in all. A query $[l, r)$ takes the largest $k$ with $2^k \le r - l$; the two blocks of length $2^k$ that start at
$l$ and end at $r$ cover the range, overlapping in the middle, and the answer is the smaller of their two minima
(Figure~\ref{fig:sparse}). That is $\Oh{1}$ per query, with $k = \lfloor \log_2 (r - l) \rfloor$ from \cc{__lg}.

\begin{figure}[htbp]
\centering
\begin{tikzpicture}
  \node[blbl, anchor=east] at (-0.2, 0) {$a$};
  \arrs{A}{0}{0}{5/n, 2/cool, 8/cool, 6/cool, 3/good, 9/good, 7/good, 4/good, 1/n, 6/n}
  \idx{A}{10}
  \foreach \lv/\vals/\yy in {1/{2,2,6,3,3,7,4,1,1}/-1.35, 2/{2,2,3,3,3,1,1}/-2.1, 3/{2,1,1}/-2.85}
    { \node[blbl, anchor=east] at (-0.2, \yy) {$k = \lv$};
      \arr{L\lv}{0}{\yy}{\vals} }
  \node[cellc, hot, anchor=west] at (L2-2.west) {2};
  \node[cellc, hot, anchor=west] at (L2-5.west) {3};
  \draw[line width=1.6pt, draw=fgblue, rounded corners=2pt] ($(A-2.north west)+(0.04, 0.42)$) rectangle
    ($(A-5.north east)+(-0.04, 0.12)$);
  \draw[line width=1.6pt, draw=fggreen!70!black, rounded corners=2pt] ($(A-5.north west)+(0.04, 0.82)$) rectangle
    ($(A-8.north east)+(-0.04, 0.52)$);
  \node[font=\sffamily\scriptsize, text=fgblue] at ($(A-2.north west)!0.5!(A-5.north east)+(0, 0.27)$) {$[1, 5)$: min 2};
  \node[font=\sffamily\scriptsize, text=fggreen!60!black] at ($(A-5.north west)!0.5!(A-8.north east)+(0, 0.67)$)
    {$[4, 8)$: min 3};
  \node[notel, anchor=west, text width=5.2cm] at (7.4, -1.1) {The query $[1, 8)$ has length 7, so $k = 2$: the blocks
    $[1, 5)$ and $[4, 8)$, with minima $\mathit{st}[2][1] = 2$ and $\mathit{st}[2][4] = 3$ (red), overlap in position 4. The
    answer is $\min(2, 3) = 2$.};
\end{tikzpicture}
\caption{A sparse table: row $k$ holds the minima of all blocks of length $2^k$. Any range is the union of two such blocks.}
\label{fig:sparse}
\end{figure}

\begin{problem}{Static Range Minimum Queries (CSES)}
Given $n \le 2 \cdot 10^5$ numbers and $q \le 2 \cdot 10^5$ queries, report the minimum of each range $[a, b]$.
\end{problem}

\begin{code}
void solve() {  // SparseTable (cpsparse) is pasted above solve()
    int n, q;
    cin >> n >> q;
    vi a(n);
    for (ll& x : a) cin >> x;
    SparseTable<ll> st(a);
    while (q--) {
        int l, r;
        cin >> l >> r;
        print(st.query(l - 1, r));  // positions [l - 1, r), 0-based
    }
}
\end{code}

\usetemplate{Sparse table}{Ctrl+Alt+T, R}{cpsparse}

\begin{pitfall}[Only for idempotent operations]
The two blocks overlap, so the elements in the middle are counted twice. For min, max, gcd, AND and OR that changes nothing; for a
sum, a product or a count it is wrong. Use prefix sums for invertible operations and the disjoint sparse table below for the
rest.
\end{pitfall}

\paragraph{LCA in $\Oh{1}$.} Writing down the vertices of a tree in the order a DFS passes through them (the full Euler tour of
Section~\ref{sec:rooted}, $2n - 1$ entries), the LCA of $u$ and $v$ is the shallowest vertex between their first occurrences. With a
sparse table over the depths, storing the vertex together with its depth, every LCA query becomes one range minimum: $\Oh{n \log n}$
preprocessing and $\Oh{1}$ per query, the fastest LCA for very many queries.

% ======================================================================================================================
\section[Disjoint sparse table]{Disjoint sparse table \hfill \plat}\label{sec:disjointsparse}

For an operation that is associative but neither idempotent nor invertible (products of matrices, compositions of functions,
knapsack-like merges), the two blocks of a query must not overlap. The \defn{disjoint sparse table} arranges the precomputed values
so that every range is the union of exactly two disjoint pieces. At level $h$, cut the array into blocks of length $2^h$; inside every
block, store for each position the aggregate from that position to the block's middle (a suffix of the left half) or from the
middle to that position (a prefix of the right half). For a query $[l, r]$ with $l < r$, the level is the highest bit where $l$ and
$r$ differ, plus one: at that level $l$ and $r$ fall into the same block but different halves, so the answer is the suffix stored at
$l$ combined with the prefix stored at $r$ (Figure~\ref{fig:disjoint}). Preprocessing is $\Oh{n \log n}$ operations and every query
exactly one.

\begin{figure}[htbp]
\centering
\begin{tikzpicture}
  \node[blbl, anchor=east] at (-0.2, 0) {$a$};
  \arrs{A}{0}{0}{3/n, 1/n, 4/n, 1/n, 5/n, 9/n, 2/n, 6/n}
  \idx{A}{8}
  \foreach \lv/\vals/\yy in {3/{9,6,5,1,5,14,16,22}/-1.35, 2/{4,1,4,5,14,9,2,8}/-2.1, 1/{3,1,4,1,5,9,2,6}/-2.85}
    { \node[blbl, anchor=east] at (-0.2, \yy) {$h = \lv$};
      \arr{H\lv}{0}{\yy}{\vals} }
  \foreach \lv/\cells in {3/{4}, 2/{2, 6}, 1/{1, 3, 5, 7}}
    \foreach \cc in \cells
      \draw[line width=1.8pt, draw=fgverm] ($(H\lv-\cc.north east)+(0, 0.05)$) -- ($(H\lv-\cc.south east)+(0, -0.05)$);
  \node[cellc, hot, anchor=west] at (H3-2.west) {6};
  \node[cellc, hot, anchor=west] at (H3-7.west) {16};
  \node[notel, anchor=west, text width=5.0cm] at (6.0, -1.6) {Sums of $3, 1, 4, 1, 5, 9, 2, 6$. Red bars are the block
    middles of each level; left of a middle a cell holds the sum from it to the middle, right of a middle the sum from the
    middle to it. The query $[1, 6]$: $1 \oplus 6 = 111_2$, so $h = 3$, and the answer is $6 + 16 = 22$ (red cells), with no
    overlap.};
\end{tikzpicture}
\caption{A disjoint sparse table. Every range is the suffix of one half-block plus the prefix of the next, at the level where the
two ends first fall into the same block.}
\label{fig:disjoint}
\end{figure}

\begin{problem}{Composing linear functions}
There are $n$ functions $f_i(x) = a_i x + b_i \pmod{10^9 + 7}$. Each of $q$ queries gives $l$, $r$ and $x$ and asks for $f_r(\cdots
f_{l+1}(f_l(x)) \cdots)$ ($n, q \le 2 \cdot 10^5$).
\solution Composition is associative: a pair $(a, b)$ followed by $(c, d)$ is $(ca, cb + d)$. It is not idempotent (applying a
function twice is not the same as once) and not invertible when some $a_i = 0$, so neither prefix sums nor a sparse table work, and
the disjoint sparse table answers each query with one composition.
\end{problem}

\begin{code}
void solve() {  // DisjointSparse (cpdisjointsparse) is pasted above solve()
    int n, q;
    cin >> n >> q;
    vector<pll> f(n);  // f_i(x) = a x + b (mod p)
    for (auto& [a, b] : f) cin >> a >> b;
    // (F then G)(x) = G(F(x)) = Ga (Fa x + Fb) + Gb: associative, but
    // neither idempotent nor invertible, so a plain sparse table fails
    auto then = [](pll F, pll G) {
        ll a = G.first * F.first % MOD;
        ll b = (G.first * F.second + G.second) % MOD;
        return pll{a, b};
    };
    DisjointSparse<pll, decltype(then)> t(f, then);
    while (q--) {
        int l, r;
        ll x;
        cin >> l >> r >> x;
        auto [a, b] = t.query(l - 1, r);  // f_l first, f_r last
        print((a * x + b) % MOD);
    }
}
\end{code}

\usetemplate{Disjoint sparse table}{Ctrl+Alt+U, T}{cpdisjointsparse}

\begin{trick}[The offline version: divide and conquer]
When all queries are known in advance, the same idea needs only $\Oh{n}$ memory: split the array at its middle $m$, compute the
suffix aggregates of the left half and the prefix aggregates of the right half, answer every query that crosses $m$ with one
combination, and recurse on the queries that lie inside a half. It also works when the ``aggregate'' is a whole DP table, such as a
bitset knapsack: ``can a subset of $a_l, \dots, a_r$ sum to $s$?'' is answered by combining the knapsacks of a suffix and a prefix,
$\Oh{(n \log n + q) W / 64}$ in all.
\end{trick}

% ======================================================================================================================
\section[Cartesian trees]{Cartesian trees \hfill \plat}\label{sec:cartesian}

The \defn{Cartesian tree} of an array has the minimum as its root (the leftmost one, if it repeats), the Cartesian tree of the part
to its left as the left subtree and the Cartesian tree of the part to its right as the right subtree (Figure~\ref{fig:cartesian}).
It is a binary search tree by position and a heap by value, and it has two useful properties.
\begin{itemize}
  \item The minimum of $a_l, \dots, a_r$ is the LCA of $l$ and $r$ in the tree; range minimum and LCA are the same problem.
  \item The subtree of $i$ is exactly the widest range in which $a_i$ is the minimum (the leftmost minimum, on ties). So $a_i$ is the
        minimum of the $(\mathit{left}_i + 1)(\mathit{right}_i + 1)$ subarrays inside that range that contain $i$, where
        $\mathit{left}_i$ and $\mathit{right}_i$ are the sizes of its subtrees.
\end{itemize}

\begin{figure}[htbp]
\centering
\begin{tikzpicture}
  \tikzset{ct/.style={vtx, minimum size=7mm}}
  \foreach \ii/\val/\dd/\sy in {0/3/1/n, 1/1/0/hot, 2/4/2/n, 3/1/1/warn, 4/5/3/n, 5/9/4/n, 6/2/2/cool, 7/6/3/n}
    \node[ct, \sy] (C\ii) at (0.9*\ii, -0.85*\dd) {\val};
  \foreach \ua/\ub in {1/0, 1/3, 3/2, 3/6, 6/4, 4/5, 6/7} \draw[ed] (C\ua) -- (C\ub);
  \node[blbl, anchor=east] at (-0.6, -4.4) {$a$};
  \foreach \val [count=\ii from 0] in {3, 1, 4, 1, 5, 9, 2, 6} {
    \node[cellc] (A\ii) at (0.9*\ii, -4.4) {\val};
    \node[lbl] at (0.9*\ii, -4.98) {\ii};
  }
  \draw[mbrace] ($(A4.south west)+(0, -0.4)$) -- node[below=6pt, lbl] {subtree of the 2: $[4, 7]$} ($(A7.south east)+(0, -0.4)$);
  \node[notel, anchor=west, text width=5.0cm] at (7.6, -1.6) {The Cartesian tree of $3, 1, 4, 1, 5, 9, 2, 6$, each vertex drawn
    above its position. The root is the first 1 (red), the second 1 (orange) is its right child. The 2 at position 6 (blue)
    rules positions 4 to 7: it is the minimum of $(2 + 1)(1 + 1) = 6$ subarrays.};
\end{tikzpicture}
\caption{A Cartesian tree is a heap by value and a binary search tree by position: in-order is array order.}
\label{fig:cartesian}
\end{figure}

\paragraph{Linear-time construction.} Scan the array from left to right and keep the \emph{right spine} of the tree built so far, a
stack of indices with increasing values. A new element $a_i$ pops every spine element larger than it; the last element popped becomes
its left child, and $a_i$ becomes the right child of the new top of the stack. Every index is pushed and popped once: $\Oh{n}$. It is
the same monotonic stack as for ``nearest smaller element'' (Section~\ref{sec:monostack-apps}), and the parent of $i$ is the larger
of its two nearest smaller neighbours.

\begin{problem}{Sum of subarray minimums, again}
Compute the sum, modulo $10^9 + 7$, of the minimum of every subarray of $n \le 2 \cdot 10^5$ numbers.
\solution Section~\ref{sec:contribution} solved this with two monotonic stacks; the Cartesian tree gives the same count as subtree
sizes: $\sum_i a_i (\mathit{left}_i + 1)(\mathit{right}_i + 1)$. Ties need no special care, because each subarray is charged to the LCA
of its ends, which is one vertex.
\end{problem}

\begin{code}
void solve() {  // CartesianTree (cpcartesian) is pasted above solve()
    int n;
    cin >> n;
    vi a(n);
    for (ll& x : a) cin >> x;
    CartesianTree ct(a);
    // the subtree of i is the widest range in which a[i] is the minimum
    // (the leftmost one, on ties): a[i] is the minimum of the subarrays in
    // it that contain i, (left subtree + 1) * (right subtree + 1) of them
    vi order{ct.root}, sub(n, 1);
    rep(i, n) for (int c : {ct.lc[order[i]], ct.rc[order[i]]})
        if (c != -1) order.push_back(c);
    per(i, n) {  // children before parents
        int v = order[i];
        for (int c : {ct.lc[v], ct.rc[v]})
            if (c != -1) sub[v] += sub[c];
    }
    ll ans = 0;
    rep(i, n) {
        ll left = ct.lc[i] == -1 ? 1 : sub[ct.lc[i]] + 1;
        ll right = ct.rc[i] == -1 ? 1 : sub[ct.rc[i]] + 1;
        ans = (ans + a[i] % MOD * (left * right % MOD)) % MOD;
    }
    print(ans);
}
\end{code}

\usetemplate{Cartesian tree}{Ctrl+Alt+U, C}{cpcartesian}

\begin{recognize}
``For every element, the range where it is the minimum (or maximum)'', ``largest rectangle in a histogram'' (each bar's rectangle is
its subtree), sums or counts over all subarrays by their minimum, and divide-and-conquer that always splits at the minimum: the
recursion tree of that divide and conquer \emph{is} the Cartesian tree, so it can be built in $\Oh{n}$ instead of searching for each
minimum.
\end{recognize}

% ======================================================================================================================
\section[Square-root decomposition]{Square-root decomposition \hfill \plat}\label{sec:sqrt}

Split the array into blocks of size $B$ and keep a summary of every block. A query covers some whole blocks, at most $n / B$ of them,
and at most $2B$ loose elements at its two ends; an update changes one element and rebuilds one summary (Figure~\ref{fig:sqrt}). With
$B \approx \sqrt n$ every operation costs $\Oh{\sqrt n}$. It is slower than a segment tree, but the summaries can be anything:
sorted copies of the block, frequency tables, or answers to questions that do not combine across blocks.

\begin{figure}[htbp]
\centering
\begin{tikzpicture}
  \foreach \bb in {0, ..., 3} {
    \foreach \ii in {0, ..., 3} {
      \pgfmathtruncatemacro{\xx}{\bb * 4 + \ii}
      \node[cellc] (Q\xx) at (\xx * 0.62 + \bb * 0.22, 0) {};
    }
    \node[seg, minimum width=2.35cm] (S\bb) at (\bb * 2.70 + 0.93, 0.9) {block \bb};
  }
  \foreach \xx in {3, 12, 13} \node[cellc, hot] at (Q\xx) {};
  \foreach \bb in {1, 2} \node[seg, minimum width=2.35cm, fill=fggreen!20, draw=fggreen!70!black] at (S\bb) {block \bb};
  \draw[mbrace] ($(Q3.south west)+(0, -0.12)$) -- node[below=6pt, lbl] {query $[3, 14)$} ($(Q13.south east)+(0, -0.12)$);
  \node[notel, anchor=west, text width=4.1cm] at (11.3, 0.45) {The query uses the summaries of the whole blocks 1 and 2
    (green) and looks at the three loose elements (red) one by one.};
\end{tikzpicture}
\caption{Square-root decomposition: at most $2B$ loose elements and $n / B$ whole blocks per query.}
\label{fig:sqrt}
\end{figure}

\begin{problem}{Counting small values with updates}
An array of $n \le 2 \cdot 10^5$ numbers receives $q \le 2 \cdot 10^5$ operations: assign a value to one element, or count the elements
$\le x$ in a range $[l, r]$.
\solution Keep every block sorted. An assignment replaces one value in its sorted block, $\Oh{B}$; a count binary-searches each whole
block and checks the loose elements, $\Oh{(n / B) \log B + B}$. Balancing gives $B \approx \sqrt{n \log n}$, a few hundred; the program
uses 512 and takes 0.46\,s for $n = q = 2 \cdot 10^5$. (A merge-sort tree or a wavelet tree answers the counts faster but cannot take the updates as easily.)
\end{problem}

\begin{code}
void solve() {
    int n, q;
    cin >> n >> q;
    vi a(n);
    for (ll& x : a) cin >> x;
    const int B = 512;  // the block size, about sqrt(n log n)
    int nb = (n + B - 1) / B;
    vvi blk(nb);  // the values of each block, kept sorted
    rep(i, n) blk[i / B].push_back(a[i]);
    for (auto& v : blk) sort(all(v));
    while (q--) {
        int type;
        cin >> type;
        if (type == 1) {  // a[i] = x: replace one value in its block
            int i;
            ll x;
            cin >> i >> x;
            i--;
            auto& v = blk[i / B];
            v.erase(lower_bound(all(v), a[i]));
            v.insert(upper_bound(all(v), x), x);
            a[i] = x;
            continue;
        }
        int l, r;
        ll x;
        cin >> l >> r >> x;
        l--;  // count the elements <= x in [l, r)
        ll cnt = 0;
        while (l < r && l % B) cnt += a[l++] <= x;  // the loose head
        for (; l + B <= r; l += B) {  // whole blocks: one binary search each
            auto& v = blk[l / B];
            cnt += upper_bound(all(v), x) - v.begin();
        }
        while (l < r) cnt += a[l++] <= x;  // the loose tail
        print(cnt);
    }
}
\end{code}

\begin{trick}[Rebuild every $\sqrt q$ operations]
When updates are hard to apply to a static structure but easy to account for one by one, keep a \emph{buffer} of the pending updates.
Every query asks the static structure and then corrects its answer for each buffered update; when the buffer reaches about $\sqrt q$
entries, rebuild the structure from scratch with all of them applied. With rebuilds in $\Oh{n}$ and corrections in $\Oh{1}$, the total is
$\Oh{(n + q)\sqrt q}$, and the structure itself never has to support updates.
\end{trick}

\begin{trick}[Heavy and light, beyond arrays]
The square-root threshold balances work in many other settings. In a graph with $m$ edges, at most $2\sqrt m$ vertices have degree
$\ge \sqrt m$; call them heavy. ``Precompute an answer for every heavy vertex, iterate over the neighbours of every light one'' keeps
every query $\Oh{\sqrt m}$: triangle counting, sums over neighbours under updates, and set intersections all use it. In strings, there
are at most $\Oh{\sqrt n}$ distinct pattern lengths when the total length is $n$, which bounds the work of hashing every length.
\end{trick}

% ======================================================================================================================
\section[Mo's algorithm]{Mo's algorithm \hfill \plat}\label{sec:mo}

Some statistics have no useful summary for a block, but are easy to update when the range grows or shrinks by one element: the
number of distinct values, the number of pairs of equal values, the most frequent value. \defn{Mo's algorithm} answers many such
queries offline by reordering them so that a window $[L, R)$ moving from query to query travels little. Sort the queries by the block
of their left end ($B \approx n / \sqrt q$) and, within a block, by their right end, alternating the direction from block to block;
then move the window step by step, calling \emph{add} and \emph{remove} (Figure~\ref{fig:mo}).

\emph{Why $\Oh{(n + q)\sqrt q}$ moves.} Within one block, the right end only moves in one direction, $\Oh{n}$ per block, and there are
$n / B = \sqrt q$ blocks: $\Oh{n \sqrt q}$. The left end moves at most $2B$ per query: $\Oh{qB} = \Oh{n \sqrt q}$. So the window moves
$\Oh{n \sqrt q}$ steps in all, about $10^8$ cheap steps for $n = q = 2 \cdot 10^5$.

\begin{figure}[htbp]
\centering
\begin{tikzpicture}
\begin{axis}[cpplot, width=8.6cm, height=6.6cm, xmin=-0.5, xmax=16.5, ymin=0, ymax=17, xlabel={left end $l$}, ylabel={right end $r$},
  xtick={0, 4, 8, 12, 16}, ytick={0, 4, 8, 12, 16}]
  \addplot[fgblue, line width=1pt, mark=*, mark size=2pt, ->] coordinates {(2,3) (1,5) (0,9) (3,15) (7,12) (5,7) (4,6) (8,12)
    (11,14) (12,16) (12,15)};
  \draw[dashed, fgline] (axis cs:3.5, 0) -- (axis cs:3.5, 17);
  \draw[dashed, fgline] (axis cs:7.5, 0) -- (axis cs:7.5, 17);
  \draw[dashed, fgline] (axis cs:11.5, 0) -- (axis cs:11.5, 17);
\end{axis}
\end{tikzpicture}
\hspace{4mm}
\begin{minipage}[b]{5.0cm}\sffamily\footnotesize\color{fgink!80}\raggedright
Twelve queries on 16 elements, in the order Mo visits them; the dashed lines separate the blocks of the left end. In the first
and third blocks $r$ increases, in the second and fourth it decreases. The window moves 57 steps in total, against 164 in the
input order.\par\vspace{1.2cm}
\end{minipage}
\caption{Mo's order: left ends grouped in blocks, right ends sweeping up and down like a snake.}
\label{fig:mo}
\end{figure}

\begin{problem}{Distinct Values Queries (CSES)}
Report the number of distinct values in each of $q \le 2 \cdot 10^5$ ranges of an array of $n \le 2 \cdot 10^5$ numbers.
\solution Compress the values to $0..n-1$ and keep a count per value: adding an element increases the number of distinct values when
its count goes from 0 to 1, removing one decreases it when its count drops to 0.
\end{problem}

\begin{code}
void solve() {  // Mo's algorithm (cpmo) is pasted above solve()
    int n, q;
    cin >> n >> q;
    vi a(n);
    for (ll& x : a) cin >> x;
    vi vals = a;  // compress the values to 0, 1, ..., n - 1
    sort(all(vals));
    vals.erase(unique(all(vals)), vals.end());
    for (ll& x : a) x = lower_bound(all(vals), x) - vals.begin();
    vector<pii> qs(q);
    for (auto& [l, r] : qs) cin >> l >> r, l--;  // [l, r)
    vi cnt(n, 0), ans(q);
    int distinct = 0;
    mo(n, qs, [&](int i) { distinct += cnt[a[i]]++ == 0; },
       [&](int i) { distinct -= --cnt[a[i]] == 0; },
       [&](int qi) { ans[qi] = distinct; });
    for (ll x : ans) print(x);
}
\end{code}

\usetemplate{Mo's algorithm}{Ctrl+Alt+Y, O}{cpmo}

This particular problem also has an $\Oh{(n + q) \log n}$ offline solution with a Fenwick tree, sweeping the right end and keeping
only the last occurrence of every value (Section~\ref{sec:sweepqueries}); Mo's algorithm is the general tool for the statistics that have
no such trick. With $n = q = 2 \cdot 10^5$ the program runs in 0.20\,s.

\begin{trick}[Making Mo faster]
\begin{itemize}
  \item \textbf{Snake order}: alternate the direction of $r$ between blocks (the template does), which saves the return trip of the
        right end, often 30--40\,\%.
  \item \textbf{Hilbert order}: sort the queries along a Hilbert curve through the $(l, r)$ plane instead; consecutive queries are
        then close in both coordinates, which is usually faster again.
  \item \textbf{Rollback Mo}: when adding is easy but removing is not (a maximum, a DSU), keep the left end at the end of its block,
        only add to the right, handle the part inside the block by adding and then undoing.
\end{itemize}
\end{trick}

\subsection{Mo with updates}

With point updates mixed into the queries, add time as a third coordinate: a query is $(l, r, t)$, the range after the first $t$
updates. Applying or undoing an update when the window passes through it is a swap of the stored value, plus a remove and an add if
the position lies inside the window. Sorting by (block of $l$, block of $r$, $t$) with blocks of $n^{2/3}$ gives $\Oh{n^{5/3}}$ moves,
about $2 \cdot 10^8$ for $n = 10^5$: heavy, but often the only way. For distinct values with $n = 10^5$ elements and $10^5$
operations, half of them updates, the program below takes 0.41\,s.

\begin{code}
void solve() {  // Mo with updates (cpmoupdates) is pasted above solve()
    int n, q;
    cin >> n >> q;
    vi a(n);
    for (ll& x : a) cin >> x;
    vector<pll> upd;  // (position, new value)
    vector<array<ll, 3>> qs;  // (l, r, updates before it), the range [l, r)
    vi vals = a;
    rep(i, q) {
        int type;
        ll x, y;
        cin >> type >> x >> y;
        if (type == 1) upd.push_back({x - 1, y}), vals.push_back(y);
        else qs.push_back({x - 1, y, sz(upd)});
    }
    sort(all(vals));
    vals.erase(unique(all(vals)), vals.end());
    auto id = [&](ll v) { return lower_bound(all(vals), v) - vals.begin(); };
    for (ll& x : a) x = id(x);
    for (auto& [p, v] : upd) v = id(v);
    vi cnt(sz(vals), 0), ans(sz(qs));
    int distinct = 0;
    auto add = [&](int i) { distinct += cnt[a[i]]++ == 0; };
    auto del = [&](int i) { distinct -= --cnt[a[i]] == 0; };
    auto apply = [&](int u, int L, int R) {  // toggle update u
        auto& [p, v] = upd[u];
        bool inside = L <= p && p < R;
        if (inside) del(p);
        swap(a[p], v);  // applying and undoing are the same swap
        if (inside) add(p);
    };
    mo_with_updates(n, qs, add, del, apply, [&](int i) { ans[i] = distinct; });
    for (ll x : ans) print(x);
}
\end{code}

\usetemplate{Mo with updates}{Ctrl+Alt+U, Shift+O}{cpmoupdates}

\subsection{Mo on trees}

Questions about the \emph{path} between two vertices (``how many distinct colours on the path from $u$ to $v$'') become range questions
through the Euler tour in which every vertex appears twice, once when the DFS enters it and once when it leaves
(Figure~\ref{fig:motree}). For $u$ entered before $v$: if $u$ is an ancestor of $v$, the path is the range from $u$'s entry to $v$'s
entry; otherwise it is the range from $u$'s exit to $v$'s entry, plus the LCA, which that range misses. In both cases the vertices of
the range that appear \emph{exactly once} are the path: a vertex whose whole subtree lies in the range appears twice and cancels.
So Mo's algorithm runs on the sequence with a ``toggle'' instead of add and remove; for distinct colours on paths with $n = q = 2
\cdot 10^5$ the program below takes 0.62\,s.

\begin{figure}[htbp]
\centering
\begin{tikzpicture}
  \foreach \nm/\px/\py/\sy in {0/1.5/0/warn, 1/0.5/-1.0/good, 2/2.5/-1.0/good, 3/0/-2.0/n, 4/1.0/-2.0/good,
      5/2.0/-2.0/good, 6/3.0/-2.0/n}
    \node[vtx, \sy] (T\nm) at (\px, \py) {\nm};
  \draw[ed] (T0) -- (T1) (T0) -- (T2) (T1) -- (T3) (T1) -- (T4) (T2) -- (T5) (T2) -- (T6);
  \node[blbl, anchor=east] at (5.3, -0.5) {Euler tour};
  \arrs{E}{5.45}{-0.5}{0/n, 1/n, 3/n, 3/n, 4/n, 4/good, 1/good, 2/good, 5/good, 5/n, 6/n, 6/n, 2/n, 0/n}
  \idx{E}{14}
  \cellframe{E-6}{E-9}{fgverm}
  \node[notel, anchor=north west, text width=9.6cm] at (5.45, -1.25) {The path from 4 to 5: 4 is entered first and is not an
    ancestor of 5, so take positions $5..8$, from the exit of 4 to the entry of 5. There, 4, 1, 2 and 5 each appear once: the path
    minus its LCA 0 (orange), which is added separately. Inside a range, a vertex seen twice is a whole subtree that the path
    does not touch.};
\end{tikzpicture}
\caption{Mo on trees: a path is a range of the doubled Euler tour, toggling vertices, plus possibly the LCA.}
\label{fig:motree}
\end{figure}

\begin{code}
void solve() {  // Mo on trees (cpmotree) is pasted above solve()
    int n, q;
    cin >> n >> q;
    vi c(n);
    for (ll& x : c) cin >> x;
    vvi adj(n);
    rep(i, n - 1) {
        int a, b;
        cin >> a >> b;
        a--, b--;
        adj[a].push_back(b);
        adj[b].push_back(a);
    }
    vi vals = c;  // compress the colours
    sort(all(vals));
    vals.erase(unique(all(vals)), vals.end());
    for (ll& x : c) x = lower_bound(all(vals), x) - vals.begin();
    vector<pii> qs(q);
    for (auto& [u, v] : qs) cin >> u >> v, u--, v--;
    vi cnt(n, 0), ans(q);
    int distinct = 0;
    mo_on_tree(adj, 0, qs, [&](int v) { distinct += cnt[c[v]]++ == 0; },
               [&](int v) { distinct -= --cnt[c[v]] == 0; },
               [&](int i) { ans[i] = distinct; });
    for (ll x : ans) print(x);
}
\end{code}

\usetemplate{Mo on trees}{Ctrl+Alt+U, O}{cpmotree}

\begin{recognize}
Many range (or path) queries, no updates or few, a statistic that is easy to maintain when one element enters or leaves but does not
combine across blocks (distinct counts, the mode, pairs of equal values, inversions with a Fenwick tree), and $n, q \le 2 \cdot 10^5$.
If the queries must be answered online, Mo's algorithm does not apply; persistent structures (Chapter~\ref{ch:segtree}) usually do.
\end{recognize}

\begin{summary}[Chapter summary]
\begin{itemize}
  \item Idempotent operations: a sparse table, two overlapping power-of-two blocks per query, $\Oh{1}$. LCA is a range minimum over
        the Euler tour.
  \item Any associative operation: a disjoint sparse table, one suffix and one prefix per query, $\Oh{1}$; offline, divide and
        conquer does the same in linear memory.
  \item The Cartesian tree, built with one monotonic stack, makes range minima LCAs and subtrees ``ranges where $a_i$ is the
        minimum''.
  \item Square-root decomposition handles summaries that nothing else supports, and periodic rebuilds replace hard updates.
  \item Mo's algorithm reorders offline queries so that a window moves $\Oh{n \sqrt q}$ steps; time as a third coordinate handles
        updates, and the doubled Euler tour handles paths.
\end{itemize}
\end{summary}

\begin{practice}
\begin{pplist}
  \item Static Range Sum Queries
  \item Static Range Minimum Queries
  \item Range Xor Queries
  \item Distinct Values Queries
  \item Advertisement
  \item Subarray Distinct Values
\end{pplist}
\end{practice}
